%============================================================
\chapter{Implementierung}
\label{chap:implementierung}
%============================================================

\section{ROS\,2-Knotenarchitektur}
\label{sec:ros_nodes}

Die gesamte Fahrzeuglogik ist in einem einzigen ROS\,2-Knoten
(\texttt{ftg\_node}) implementiert.
\Cref{fig:ros_graph} zeigt den vereinfachten Kommunikationsgraphen.

\begin{figure}[H]
  \centering
  \fbox{\parbox{0.85\textwidth}{\centering\vspace{2.5cm}
        Abbildung: ROS\,2-Knoten-Graph\\[0.5em]
        \texttt{/scan} $\rightarrow$ \texttt{ftg\_node}
        $\rightarrow$ \texttt{/drive}\\[1em]
        (Bild in \texttt{images/ros\_graph.png} hinterlegen)
        \vspace{2.5cm}}}
  \caption{ROS\,2-Kommunikationsgraph des FTG-Systems}
  \label{fig:ros_graph}
\end{figure}

\begin{table}[H]
  \centering
  \caption{Topics des \texttt{ftg\_node}}
  \label{tab:topics}
  \begin{tabular}{@{}llll@{}}
    \toprule
    \textbf{Topic}  & \textbf{Richtung} & \textbf{Typ}
                    & \textbf{Beschreibung} \\
    \midrule
    \texttt{/scan}  & Subscribe  & \texttt{sensor\_msgs/LaserScan}
                    & Rohe LiDAR-Daten \\
    \texttt{/drive} & Publish    & \texttt{ackermann\_msgs/AckermannDriveStamped}
                    & Lenkwinkel und Geschwindigkeit \\
    \bottomrule
  \end{tabular}
\end{table}

%------------------------------------------------------------
\section{Quellcode}
\label{sec:source_code}
%------------------------------------------------------------

\subsection{Hauptknoten (\texttt{ftg\_node.py})}
\label{subsec:ftg_node}

\begin{lstlisting}[language=Python3,
                   caption={FTG-Hauptknoten (vereinfacht)},
                   label={lst:ftg_node}]
#!/usr/bin/env python3
"""
Follow-the-Gap Node for the Mxck F1TENTH vehicle.
Subscribes to /scan and publishes Ackermann drive commands to /drive.
"""

import math
import numpy as np
import rclpy
from rclpy.node import Node
from sensor_msgs.msg import LaserScan
from ackermann_msgs.msg import AckermannDriveStamped


class FollowTheGapNode(Node):
    """Reactive obstacle avoidance using the Follow-the-Gap algorithm."""

    # ------------------------------------------------------------------ #
    #  Tuning parameters                                                   #
    # ------------------------------------------------------------------ #
    BUBBLE_RADIUS   = 0.35   # [m] safety bubble radius
    MIN_RANGE       = 0.10   # [m] ignore points closer than this
    MAX_RANGE       = 3.00   # [m] cap sensor range
    ANGLE_FAST      = 10.0   # [deg] below -> max speed
    ANGLE_SLOW      = 20.0   # [deg] above -> min speed
    SPEED_MAX       = 1.5    # [m/s]
    SPEED_MID       = 1.0    # [m/s]
    SPEED_MIN       = 0.5    # [m/s]

    def __init__(self) -> None:
        super().__init__("ftg_node")
        self.sub = self.create_subscription(
            LaserScan, "/scan", self._scan_callback, 10
        )
        self.pub = self.create_publisher(
            AckermannDriveStamped, "/drive", 10
        )
        self.get_logger().info("Follow-the-Gap node started.")

    # ------------------------------------------------------------------ #
    #  Callback                                                            #
    # ------------------------------------------------------------------ #
    def _scan_callback(self, msg: LaserScan) -> None:
        ranges = np.array(msg.ranges, dtype=np.float32)
        angles = np.linspace(
            msg.angle_min, msg.angle_max, len(ranges)
        )

        # 1. Preprocess ranges
        ranges = self._preprocess(ranges)

        # 2. Safety bubble
        ranges = self._apply_bubble(ranges)

        # 3. Find best gap
        start, end = self._find_best_gap(ranges)

        # 4. Steering target
        goal_idx = (start + end) // 2
        steering = float(angles[goal_idx])

        # 5. Velocity
        speed = self._speed_from_angle(steering)

        self._publish(steering, speed)

    # ------------------------------------------------------------------ #
    #  Algorithm steps                                                     #
    # ------------------------------------------------------------------ #
    def _preprocess(self, ranges: np.ndarray) -> np.ndarray:
        ranges = np.clip(ranges, self.MIN_RANGE, self.MAX_RANGE)
        ranges = np.nan_to_num(ranges, nan=self.MAX_RANGE,
                               posinf=self.MAX_RANGE)
        return ranges

    def _apply_bubble(self, ranges: np.ndarray) -> np.ndarray:
        closest = int(np.argmin(ranges))
        r_close = ranges[closest]
        if r_close < self.BUBBLE_RADIUS:
            half = int(math.ceil(
                math.atan2(self.BUBBLE_RADIUS, r_close)
                / (2 * math.pi / len(ranges))
            ))
            lo = max(0, closest - half)
            hi = min(len(ranges) - 1, closest + half)
            ranges[lo : hi + 1] = 0.0
        return ranges

    def _find_best_gap(self, ranges: np.ndarray):
        """Return (start, end) indices of the widest free gap."""
        free = ranges > self.MIN_RANGE
        best = (-1, -1)
        best_len = 0
        i = 0
        while i < len(free):
            if free[i]:
                j = i
                while j < len(free) and free[j]:
                    j += 1
                if j - i > best_len:
                    best_len = j - i
                    best = (i, j - 1)
                i = j
            else:
                i += 1
        if best == (-1, -1):
            mid = len(ranges) // 2
            return mid, mid
        return best

    def _speed_from_angle(self, angle_rad: float) -> float:
        deg = abs(math.degrees(angle_rad))
        if deg <= self.ANGLE_FAST:
            return self.SPEED_MAX
        if deg <= self.ANGLE_SLOW:
            return self.SPEED_MID
        return self.SPEED_MIN

    # ------------------------------------------------------------------ #
    #  Publisher helper                                                    #
    # ------------------------------------------------------------------ #
    def _publish(self, steering: float, speed: float) -> None:
        msg = AckermannDriveStamped()
        msg.drive.steering_angle = steering
        msg.drive.speed = speed
        self.pub.publish(msg)


def main(args=None):
    rclpy.init(args=args)
    node = FollowTheGapNode()
    rclpy.spin(node)
    node.destroy_node()
    rclpy.shutdown()


if __name__ == "__main__":
    main()
\end{lstlisting}

%------------------------------------------------------------
\section{Parameter-Übersicht}
\label{sec:params}
%------------------------------------------------------------

\begin{table}[H]
  \centering
  \caption{Tuning-Parameter des FTG-Algorithmus}
  \label{tab:parameter}
  \begin{tabular}{@{}lrrl@{}}
    \toprule
    \textbf{Parameter}        & \textbf{Symbol}
      & \textbf{Wert}         & \textbf{Beschreibung} \\
    \midrule
    Blasenradius    & $r_{\text{bubble}}$ & $0{,}35\,\text{m}$
      & Sicherheitsabstand zur Fahrzeugbreite \\
    Min-Abstand     & $r_{\min}$          & $0{,}10\,\text{m}$
      & Ignoriert zu nahe Punkte \\
    Max-Abstand     & $r_{\max}$          & $3{,}00\,\text{m}$
      & Kappungsgrenze für Sensor-Maximalwerte \\
    Winkel (schnell)& $\alpha_1$          & $10°$
      & Grenze für $v_{\max}$ \\
    Winkel (langsam)& $\alpha_2$          & $20°$
      & Grenze für $v_{\text{mid}}$ \\
    $v_{\max}$      & –                   & $1{,}5\,\text{m/s}$
      & Maximale Fahrgeschwindigkeit \\
    $v_{\text{mid}}$& –                   & $1{,}0\,\text{m/s}$
      & Mittlere Geschwindigkeit \\
    $v_{\min}$      & –                   & $0{,}5\,\text{m/s}$
      & Mindestgeschwindigkeit \\
    \bottomrule
  \end{tabular}
\end{table}

%------------------------------------------------------------
\section{Launch-Datei}
\label{sec:launch}
%------------------------------------------------------------

\begin{lstlisting}[language=Python3,
                   caption={Launch-Datei \texttt{ftg\_launch.py}},
                   label={lst:launch}]
from launch import LaunchDescription
from launch_ros.actions import Node


def generate_launch_description():
    return LaunchDescription([
        Node(
            package="mxck_ftg",
            executable="ftg_node",
            name="ftg_node",
            output="screen",
            parameters=[{
                "bubble_radius": 0.35,
                "min_range":     0.10,
                "max_range":     3.00,
                "speed_max":     1.5,
                "speed_mid":     1.0,
                "speed_min":     0.5,
            }],
        ),
    ])
\end{lstlisting}
